% Generated draft: requires mathematical and source review.
\documentclass[11pt]{article}
\usepackage[margin=25mm]{geometry}
\usepackage{fontspec}
\setmainfont{Latin Modern Roman}
\usepackage{amsmath,amssymb,mathtools,unicode-math}
\setmathfont{Latin Modern Math}
\usepackage{enumitem,needspace,xurl,hyperref}
\setlength{\parindent}{0pt}
\setlength{\parskip}{5pt}
\setlength{\emergencystretch}{3em}
\newcommand{\N}{\mathbb{N}}
\newcommand{\Z}{\mathbb{Z}}
\newcommand{\Q}{\mathbb{Q}}
\newcommand{\R}{\mathbb{R}}
\newcommand{\C}{\mathbb{C}}
\newcommand{\F}{\mathbb{F}}
\newcommand{\A}{\mathbb{A}}
\newcommand{\PP}{\mathbb{P}}
\newcommand{\E}{\mathbb{E}}
\newcommand{\eps}{\varepsilon}
\newcommand{\dd}{\,\mathrm{d}}
\begin{document}
\section*{UnsolvedMath ID: 20000432}
\textbf{Source problem code:} catalogue entry
\subsection{Definitions and mathematical statement}
An arithmetic scheme in the source means a scheme of finite type over \(B=\operatorname{Spec}\Z\). Flach asks for a definition of its Weil--\'etale topos, using the then-current Weil--\'etale topology, as a fibred product, and for computations of constant-sheaf integral cohomology, in particular for \(\A^1_{\Z}\) and \(\PP^1_{\Z}\).

The relevant product is a \(2\)-fibre product of topoi: given geometric morphisms \(\mathcal X\to\mathcal B\) and \(\mathcal W\to\mathcal B\), the topos \(\mathcal X\times_{\mathcal B}\mathcal W\) has the universal property of pairs of geometric morphisms equipped with an isomorphism of their composites to \(\mathcal B\). Schematically, once a base Weil--\'etale topos \(B_W\to B_{\mathrm{\acute et}}\) is fixed, the ordinary requested construction has the form
\[
 X_W=X_{\mathrm{\acute et}}\times_{B_{\mathrm{\acute et}}}B_W.
\]
The analogous Artin--Verdier compactified construction requires the corresponding compactified topoi and compatible morphisms. The question itself does not supply these base definitions or archimedean conventions. For a chosen construction, the target groups are \(H^i(X_W,\Z)\), with \(\Z\) the constant sheaf and \(i\ge0\); ordinary and compactified cohomology must be distinguished.

The supplied research assessment attributes ordinary and Artin--Verdier compactified \(2\)-fibre-product constructions to Flach--Morin (2012). Its candidate obstruction uses the structural morphism \(p:X\to B\) and a standard section \(s:B\to X\) for \(X=\A^1_{\Z}\) or \(\PP^1_{\Z}\). It claims that the induced morphisms at the actual fibre-product-topos level satisfy \(p_Ws_W=\operatorname{id}\), and hence induce a split pair
\[
 R\Gamma(B_W,\Z)\ \xrightarrow{p_W^*}\ R\Gamma(X_W,\Z)
 \ \xrightarrow{s_W^*}\ R\Gamma(B_W,\Z),\qquad
 s_W^*p_W^*=\operatorname{id}.
\]
The analogous assertion is made for the compactified construction. Thus base cohomology is a derived direct summand. In the compactified theory the report says that the base has a non-finitely-generated degree-four group, forcing the line-topos \(H^4\) to be non-finitely generated as well.

Separately, it attributes to the 2018 Flach--Morin \emph{derived} theory an unconditional projective-bundle decomposition for \(\PP^1_{\Z}\) with finite degree-four cohomology. The resulting claimed non-equivalence is specifically between that later derived complex and integral constant-sheaf cohomology of the 2012 compactified fibre-product topos. It is not an identification of their coefficients or a complete computation of either line's cohomology. The exact base group, complete derived complexes and projective-bundle formula are not supplied here. Their definitions and the compactification functoriality needed for the retraction must be checked in the cited constructions before this reported obstruction is accepted.
\subsection{Short English statement}
Define a Weil--\'etale topos for every arithmetic scheme and compute the constant-sheaf cohomology of the affine and projective lines. The supplied report describes a fibre-product construction and a section-retraction obstruction distinguishing its compactified integral cohomology from a later derived theory. It does not give a full list of the requested groups.
\subsection{Sources}
\begin{itemize}
\item Local UnsolvedMath registry, permanent ID \(20000432\), source code \texttt{catalogue entry}; AIM algebraic number theory corpus, \texttt{aim-algebraic-number-theory-notes.json}, notes[83], stored source item 1. Attribution: Ulam AI / UnsolvedMath Contributors, CC BY 4.0.
\item AIM workshop \emph{The Tate conjecture}, attributed to M. Flach: \url{https://aimath.org/WWN/tateconjecture/tateconjecture.pdf}. The supplied verification reports page 5, actual Problem 21 rather than OCR ``1''; the trailing ``5'' is a page artifact. The word ``current'' is historical source wording, not a claim about a present-day definition.
\item The research summary identifies Flach--Morin's 2012 topos construction and 2018 derived theory by authors and years only; titles, precise definitions, formulas and full bibliographic citations are absent. Original remarks and literature fields are empty.
\item Registry dated August 14, 2026: statement exact, status \texttt{partially\_solved}, classification \texttt{PARTIAL-PROGRESS}; attempt 1 labelled partial result, rigor proved, novelty confidence low. The section-retraction and non-equivalence claims are machine-generated research aids, not peer reviewed, and require independent verification before citation.
\end{itemize}
\end{document}
